\color{unidarkgreen}{\subsection{Блокировка при качаниях}}
\color{black}

\begin{enumerate}
%%%%%%%%%%%%%%%%%%%%%%%%%%%%%%%%%%%%%%%%%%%%%%%%%% ОБЩИЕ СВЕДЕНИЯ %%%%%%%%%%%%%%%%%%%%%%%%%%%%%%%%%%%%%%%%%%%%%%%%%%
\item Общие сведения

\begin{enumerate}

\item
Блокировка при качаниях (БК) предназначена для исключения ложных срабатываний ДЗ при
возникновении качаний и асинхронном ходе.

\item 
В устройстве реализовано два способа блокировки:

\begin{itemize}
\item по приращению токов прямой и обратной последовательности -- БК-I;
\item по скорости изменения величины сопротивления $dZ/dt$ -- БК-Z.
\end{itemize}

\end{enumerate}


%%%%%%%%%%%%%%%%%%%%%%%%%%%%%%%%%%%%%%%%%%%%%%%%%% БК-I %%%%%%%%%%%%%%%%%%%%%%%%%%%%%%%%%%%%%%%%%%%%%%%%%%
\item  Блокировка по приращению токов прямой и обратной последовательности (БК-I)

\begin{enumerate}

\item 
БК-I принимает сигналы от чувствительного и грубого измерительных органов (ИО), контролирующих
приращения токов прямой и обратной последовательностей. На выходе блокировки при качаниях формируется
два сигнала:
\begin{itemize}
\item \textit{БК-I-б} -- вводит в работу быстродействующие ступени ДЗ;
\item \textit{БК-I-м} -- вводит в работу медленнодействующие ступени ДЗ.
\end{itemize}

\item
Функциональная схема логики БК по приращению токов прямой и обратной последовательности (БК-I) приведена на рисунке \ref{pict_bk:BK_I}.

\medskip

\begin{figure}[H]
\centering
\includegraphics[width=0.8\textwidth,height=0.8\textheight,keepaspectratio]{\fbpath/08.01 СВ/661.1471 БК/img/BK_I.pdf}
\captionof{figure}{Функциональная схема БК-I}\label{pict_bk:BK_I}
\end{figure}

\item 
Пороги срабатывания ИО, реагирующих на несимметричные КЗ, задаются с помощью уставок \textbf{«dI2 чув»}
и \textbf{«dI2 груб»}. Пороги срабатывания ИО, реагирующих на симметричные КЗ, задаются с помощью уставок
\textbf{«dI1 чув»} и \textbf{«dI1 груб»}. Срабатывание чувствительных ИО разрешает ввод в работу быстродействующих
ступеней (сигнал \textit{«БК-I-б»}) на время, задаваемое уставкой \textbf{«Тв БК-I-б чув»}, с последующей блокировкой ввода
быстродействующих ступеней ДЗ от чувствительных ПО на время \textbf{«Тв БК-I-м»}. Если в течение времени ввода
медленнодействующих ступеней \textbf{«Тв БК-I-м»} срабатывает грубый ИО, то сигнал \textit{«БК-I-б»} повторно введется в
работу, на время задаваемое уставкой \textbf{«Тв БК-I-б груб»}. Грубый ИО предусмотрен для обеспечения
повторного пуска быстродействующих ступеней ДЗ при переходе внешнего КЗ во внутреннее. Срабатывание
чувствительного или грубого ИО действует на ввод медленнодействующих ступеней ДЗ на время, задаваемое
уставкой \textbf{«Тв БК-I-м»}. В логике БК-I предусмотрен ускоренный возврат схемы при отключении выключателя,
который вводится программным ключом \textbf{«Сброс БК-I от РПО»}.

\item
Оперативный ввод/вывод БК-I производится с помощью БУ (блоков управления функциями) \textit{«Вывод БК»}.

\item
Уставочные параметры функции БК по приращению токов прямой и обратной последовательности (БК-I) представлены в таблице~\ref{app_set:bk}.

\item
Временная диаграмма работы БК по приращению токов прямой и обратной последовательности (БК-I) приведена на рисунке \ref{pict_bk:BK_I_diagram}.

\medskip

\begin{figure}[H]
\centering
\includegraphics[width=0.8\textwidth,height=0.8\textheight,keepaspectratio]{\fbpath/08.01 СВ/661.1471 БК/img/BK_I_diagram.pdf}
\captionof{figure}{Временная диаграмма работы БК-I}\label{pict_bk:BK_I_diagram}
\end{figure}

\end{enumerate}


%%%%%%%%%%%%%%%%%%%%%%%%%%%%%%%%%%%%%%%%%%%%%%%%%% БК-Z %%%%%%%%%%%%%%%%%%%%%%%%%%%%%%%%%%%%%%%%%%%%%%%%%%
\item  Блокировка по скорости изменения величины сопротивления dZ/dt (БК-Z)

\begin{enumerate}

\item
БК-Z основан на измерении времени прохождения годографом полного сопротивления области между
внешней и внутренней характеристик реле сопротивления.

\item
Функциональная схема логики БК по скорости изменения величины сопротивления dZ/dt (БК-Z) приведена на рисунке \ref{pict_bk:BK_Z}.

\medskip

\begin{figure}[H]
\centering
\includegraphics[width=0.8\textwidth,height=0.8\textheight,keepaspectratio]{\fbpath/08.01 СВ/661.1471 БК/img/BK_Z.pdf}
\captionof{figure}{Функциональная схема БК-Z}\label{pict_bk:BK_Z}
\end{figure}

\item
В качестве внутренней области характеристики можно выбрать реле сопротивления с ненаправленной
характеристикой срабатывания второй или третьей ступени ДЗ -- выбор осуществляется программным
переключателем \textbf{«ИО БК-Z»}. Внешняя характеристика срабатывания реле сопротивления отстоит от
внутренней характеристики по оси R на величину $\Delta$R и по оси $X$ на величину $\Delta$X.

\item
Характеристики реле сопротивления БК по скорости изменения величины сопротивления dZ/dt (БК-Z) приведена на рисунке \ref{pict_bk:char_relay_res}.

\begin{figure}[H]
\centering
\includegraphics[width=0.45\textwidth,height=0.45\textheight,keepaspectratio]{\fbpath/08.01 СВ/661.1471 БК/img/char_relay_res.pdf}
\captionof{figure}{Характеристики реле сопротивления, используемые для блокировки при качаниях по скорости изменения сопротивления}\label{pict_bk:char_relay_res}
\end{figure}

\item 
Значения параметров $\Delta$R и $\Delta$X не регулируются и зависят от номинального тока устройства.
При $I_{\textsl{ном}}$= 5 А:  $\Delta R$ = $\Delta X$ = 1 Ом. При $I_{\textsl{ном}}$= 1 А: $\Delta R$ = $\Delta X$ = 5 Ом.

\item
Уставка по скорости изменения Z задается выдержкой времени \textbf{«Тзад БК-Z»} – если в течение этого
времени годограф сопротивления не пересёк зону контроля скорости изменения сопротивления, то ступени
ДЗ, вводимые от БК-Z, заблокируются. Принужденный возврат схемы БК-Z задается выдержкой времени
\textbf{«Тв БК-Z»}.

\item
При возникновении короткого замыкания вектор сопротивления скачкообразно переходит из области
нагрузки в область срабатывания. При возникновении синхронных качаний вектор сопротивления появляется
в области срабатывания и покидает её. Качания выявляются при прохождении по монотонной траектории.
Логическая цепочка БК-Z выдаёт при этом запрет на срабатывание ступеней дистанционной защиты.
Срабатывание измерительного органа $I_2$ БК-Z во время качаний приводит к быстрому возврату БК-Z, и таким
образом, делает возможным отключение от дистанционной защиты. Порог срабатывания ИО $I_2$ БК-Z
регулируется уставкой \textbf{«I2 БК-Z»}. Для БК-Z также предусмотрен ускоренный возврат схемы при отключении 
выключателя, который вводится программным ключом \textbf{«Сброс БК-Z от РПО»}.

\item
Оперативный ввод/вывод БК-Z производится с помощью БУ (блоков управления функциями) \textit{«Вывод БК»}.

\item
Уставочные параметры функции БК по скорости изменения величины сопротивления dZ/dt (БК-Z) представлены в таблице~\ref{app_set:bk}.

\end{enumerate}

\end{enumerate}